\section{Input simulations}

The challenge employs simulated galaxy catalogs derived from two
complementary N-body cosmological simulations: the Cardinal
simulations and the Flagship simulation. These synthetic datasets
provide a controlled environment where the true redshifts are known by
construction, enabling rigorous validation of photometric redshift
algorithms and systematic assessment of their performance
characteristics.

The Cardinal simulations comprise a suite of high-resolution N-body
simulations specifically designed to explore the sensitivity of
cosmological observables to variations in fundamental cosmological
parameters. The simulations employ state-of-the-art semi-analytic
models to populate dark matter halos with galaxies, incorporating
realistic prescriptions for star formation, dust attenuation, and
spectral energy distribution modeling.

The Flagship simulation represents a single, ultra-large cosmological
simulation run with fiducial cosmological parameters consistent with
current observational constraints. With a volume exceeding several
cubic gigaparsecs, the Flagship provides statistical power to probe
rare objects and the high-mass end of the galaxy population. Its
primary purpose in the photometric redshift challenge is to provide a
realistic mock catalog that captures the full complexity of galaxy
populations across cosmic time, including correlations between galaxy
properties, environmental dependencies, and the intricate
relationships between spectral features and redshift.

Together, these complementary simulation suites enable challenge
participants to test both the accuracy and the robustness of their
photometric redshift estimation methods under realistic observational
conditions.
